% ECONOMICS_PROBLEM: {"id": "G23-466", "title": "Private-market valuation, liquidity and systemic links", "jel_code": "G23", "source_id": "OP-0100"}
% !TeX program = xelatex
\documentclass[11pt,a4paper]{article}
\usepackage[margin=23mm]{geometry}
\usepackage{fontspec}
\setmainfont{Latin Modern Roman}
\usepackage{amsmath,amssymb,mathtools}
\usepackage{xcolor,enumitem,booktabs,longtable,array,xurl,hyperref}
\definecolor{muted}{HTML}{53636C}
\hypersetup{colorlinks=true,urlcolor=blue,linkcolor=blue}
\setlength{\parindent}{0pt}
\setlength{\parskip}{5pt}
\newcommand{\R}{\mathbb R}
\newcommand{\E}{\mathbb E}
\newcommand{\N}{\mathbb N}
\newcommand{\ind}{\mathbf 1}
\newcommand{\Var}{\operatorname{Var}}
\newcommand{\Cov}{\operatorname{Cov}}
\newcommand{\supp}{\operatorname{supp}}
\DeclareMathOperator*{\argmin}{arg\,min}
\DeclareMathOperator*{\argmax}{arg\,max}
\newcommand{\entrymeta}[3]{{\sffamily\small\color{muted}\textbf{\texttt{#1}}\quad|\quad #2\par #3\par}}
\newcommand{\Problemref}[1]{\texttt{#1}}
\newcommand{\JELref}[2]{\texttt{#2} (JEL \texttt{#1})}
\begin{document}
\section*{Private-market valuation, liquidity and systemic links}
\textbf{Problem ID:} \texttt{G23-466}\quad \textbf{JEL:} \texttt{G23}\par
% BEGIN ECONOMICS STATEMENT
\label{problem:OP-0100}
\label{atlas:prob:F10}

\noindent\textbf{Atlas ID:} F10; \textbf{original number:} 100; \textbf{field:} Finance. \textbf{Classification:} Editorial research formulation (theoretical/empirical); not a certified unsolved theorem.

\subsection*{Definitions and assumptions}
Let private investment $i$ have contributions $C_{it}\ge0$, distributions $D_{it}\ge0$, observed appraisals $A_{it}$ and latent economically relevant asset value $V_{it}\ge0$ over dates $0,\ldots,H$. Model $\theta$ specifies cash-flow generation, fees, leverage, exit decisions, appraisal reporting and secondary-market trading with feasible contracts. Asset pricing uses an integrable positive cumulative discount factor $M_{0t}$, normalized by $M_{00}=1$, appropriate to the investor’s risks and trading opportunities. A selected equilibrium determines exit timing and secondary prices. Reporting, survival and realized-exit selection are part of the observation mechanism rather than assumptions of complete sampling.

\subsection*{Formal research question}
Identify bounds on risk-adjusted net performance, liquidity costs and common exposures using
\[
 \mathrm{NPV}_{i,\theta}=\mathbb E_\theta\!\left[
 \sum_{t=0}^{H}M_{0t}(D_{it}-C_{it})+M_{0H}V_{iH}\right],
\]
where $V_{iH}$ excludes value already paid in the terminal distribution. Determine when cash flows, secondary transactions and financing covenants identify latent value volatility and liquidation discounts rather than only realized fund performance. Estimate systemic losses under correlated refinancing or redemption shocks, tracing shared lenders and ultimate investors. Distinguish entrepreneurial holdings, buyout funds, venture capital and private credit as separate populations.

\subsection*{Interpretation and scope}
An internal rate of return can fail to exist or have multiple roots; it is not a general risk-adjusted performance statistic. A public-market equivalent uses a specified benchmark and cash-flow timing but does not automatically remove leverage, sector and liquidity-risk differences. Appraisal smoothing alone is not a uniquely identified observation model; many latent volatility paths may generate similar reported series. Secondary sales are selected by sellers’ liquidity needs, information and contract constraints, so transaction prices need a selection argument before extrapolation. Fees and carry must use a consistent investor-level cash-flow definition. Systemic risk reflects financing and counterparty links in addition to an individual asset’s illiquidity.

\subsection*{Source audit and status}
[S1] and [S2] concern institutional private-equity fund performance. [S3], also linked source [51], concerns household entrepreneurial equity; it cannot be generalized unchanged to venture funds or private credit. These historical findings do not establish a universally persistent private-equity premium puzzle. The editorial question concerns latent valuation, liquidity and systemic identification, without certification of a currently open theorem.

\subsection*{References}

\begin{enumerate}[label={[S\arabic*]},leftmargin=*]

\item Steven N. Kaplan and Antoinette Schoar (2005). \emph{Private Equity Performance: Returns, Persistence, and Capital Flows}. Journal of Finance 60(4), 1791–1823. \url{https://doi.org/10.1111/j.1540-6261.2005.00780.x}.
\textit{Locator and role:} Publisher or author abstract / bibliographic record; Fund cash flows, performance, persistence and public-market equivalents. Provides background or a benchmark result; does not establish that the editorial question remains unresolved in 2026.

\item Robert S. Harris, Tim Jenkinson, and Steven N. Kaplan (2014). \emph{Private Equity Performance: What Do We Know?}. Journal of Finance 69(5), 1851–1882. \url{https://doi.org/10.1111/jofi.12154}.
\textit{Locator and role:} Publisher or author abstract / bibliographic record; Institutional buyout and venture-fund cash-flow performance. Provides background or a benchmark result; does not establish that the editorial question remains unresolved in 2026.

\item Tobias J. Moskowitz and Annette Vissing-Jørgensen (2002). \emph{The Returns to Entrepreneurial Investment: A Private Equity Premium Puzzle?}. American Economic Review 92(4), 745–778. \url{https://doi.org/10.1257/00028280260344452}.
\textit{Locator and role:} Publisher or author abstract / bibliographic record; Concentrated entrepreneurial equity risk and measured returns; resolves link [51]. Entrepreneurial holdings are not the same population as institutional private-equity, venture or private-credit funds.

\end{enumerate}

\subsection*{Common conventions: Source mathematical conventions}

Unless an entry states otherwise, agent and firm sets are finite. State and action spaces are measurable, strategies and outcome maps are measurable, probability laws are specified, and expectations exist for the targets under discussion. Infinite-horizon discount factors lie in $(0,1)$ and discounted objectives are finite. Norms, tolerances and complexity input sizes must be specified for computational comparisons. Constants or primitives that appear locally are inputs, not empirical facts asserted by this catalogue.

\textbf{Identification.} A statistical model $m\in\mathcal M$ implies an observable law $P_m$ and target $\theta(m)$. At law $P$, its identified set is
\[
\Theta(P;\mathcal M)=\{\theta(m):m\in\mathcal M,\ P_m=P\}.
\]
Point identification holds when this set is a singleton, or equivalently when $P_{m_1}=P_{m_2}$ implies $\theta(m_1)=\theta(m_2)$ for all compatible models. Sharp bounds mean bounds attained, or approached when the boundary is not attained, by compatible models. Writing this set is a definition, not a solution: the substantive task is to establish informative restrictions and derive or estimate the set. An unrestricted-confounding model may give only trivial bounds.

\textbf{Inference.} Identification, estimability and valid uncertainty quantification are separate questions. Fix $\alpha\in(0,1)$. A confidence procedure $C_n$ over a declared law class $\mathcal P$ has asymptotic uniform coverage at level $1-\alpha$ when
\[
\liminf_{n\to\infty}\inf_{P\in\mathcal P}\Pr_P\{\theta(P)\in C_n\}\geq1-\alpha.
\]
For set-identified targets the coverage event must be defined separately, for example coverage of the full identified set. If an entry uses identification-indexed models instead of uniquely indexed laws, the same coverage convention is applied over those models. Pointwise limits do not establish uniform validity, and asymptotic validity does not guarantee finite-sample precision. Sample sizes, dependence structures and weak-identification sequences are part of each application.

\textbf{Counterfactuals.} $Y_i(a)$ is a potential outcome under a specified intervention, including its timing, implementation and versions. It is not $Y_i$ conditional on voluntarily choosing $a$. A full intervention vector permits interference; shorthand scalar treatment notation does not silently impose no interference. Assignment, selection, attrition, anticipation and measurement must be specified. A claim about an instrument's local effect is not automatically a population policy effect.

\textbf{Economic equilibrium.} A counterfactual model includes best responses, constraints, feasibility, market clearing, state transitions and expectations consistent with the stated information. When several equilibria exist, either a declared selector is an input or the target is set-valued. Comparisons across policy regimes must use the stated selection convention. Reduced-form relationships are not presumed invariant after an equilibrium-changing intervention.

\textbf{Optimization and welfare.} Optimization is over the stated feasible set. If attainment is not established, $\sup$ or $\inf$ and $\varepsilon$-optimal choices replace an asserted maximizer or minimizer. For example, with value $V=\sup_{a\in\mathcal A}W(a)<\infty$, an $\varepsilon$-optimal set is $\{a\in\mathcal A:W(a)\geq V-\varepsilon\}$ for $\varepsilon>0$. Benefits, costs, risk and health or environmental outcomes can be aggregated only using declared units and conversion weights. Interpersonal utility weights, the treatment of fiscal transfers and foreign profits, discounting, and the behavioral welfare criterion are explicit inputs. A comparative-static derivative additionally requires existence and appropriate differentiability of the selected counterfactual.

\textbf{Sources.} Local labels [S1], [S2], and so forth restart in every question. Locators identify the relevant abstract, model, section or result at the level actually checked; a bibliographic or abstract-level verification is not represented as inspection of an entire proof. Working papers are distinguished from later publications and changing versions. Source summaries are paraphrased. All URL checks and editorial formulations refer to the audit date stated above.
% END ECONOMICS STATEMENT
\end{document}
